\documentclass[10pt]{../sigplanconf}

\usepackage[utf8x]{inputenc}
\usepackage[T1]{fontenc}

\usepackage{amsmath}
\usepackage{graphicx}
\usepackage{semantic}

\conferenceinfo{DIKU workshop on Topics in Programming Languages,}{June 2011.} 
\title{To Trim a Perfect Process Tree}
%\subtitle{Subtitle Text, if any}
\authorinfo{Martin Dybdal}
           {dybber@dybber.dk}
           {DIKU, Department of Computer Science, University of Copenhagen}
\authorinfo{Ulrik Rasmussen}
           {dolle@diku.dk}
           {DIKU, Department of Computer Science, University of Copenhagen}
\authorpermission
\copyrightdata{abstract}


\mathlig{::=}{\quad ::=\quad}
\mathlig{|->}{\mapsto}

% For proofs, As frac but does not change the font size
\newcommand{\nfrac}[2]{\frac{\displaystyle{#1}}{\displaystyle{#2}}}
\newcommand{\tagsc}[1]{\tag{\scshape #1}}

\begin{document}
\maketitle


\begin{abstract}
  The Universal Resolving Algorithm\cite{abramov2000universal} is an
  algorithm for inverse interpretation, based on the notion of a
  perfect process tree for representing the possible traces of a
  program. A process tree will often have an infinite representation,
  meaning that inverse interpretation does not always terminate.

  We investigate two methods for minimizing the size of a process tree
  and representing it in finite form, respectively.

  The Size-Change Termination Analysis (SCTA)\cite{lee2001size} can be
  used to identify programs that terminate for all inputs, and
  size-change graphs resulting from this approach may be used to cut
  branches in a process tree, minimizing its size. We are not aware of
  any uses of SCTA in the field of inverse interpretation, but in this
  report we will investigate whether this approach is feasible.

  Generation of process trees by driving is a technique that
  originates from the field of
  supercompilation\cite{sorensen1998introduction}. A positive
  supercompiler uses a technique called \emph{generalization} to
  ensure that it will eventually arrive at a process tree with a
  finite representation, which still represents all possible
  configurations of the program. To our knowledge, the same technique
  has not been applied to the field of inverse interpretation.
\end{abstract}

\keywords Inverse Interpretation, Size Change Termination Analysis,
Generalization, Universal Resolving Algorithm.

\newpage
\begin{center}
  \includegraphics[width=0.7\columnwidth]{../figures/pruning.pdf}
\end{center}

\begin{figure}\centering
  \begin{align*}
    \mathcal{Q} \ni q ::= & d^{+} \\
    \mathcal{D} \ni d ::= &\texttt{fun \textit{\rmfamily fname} $x^{*}$ = $t$} \\
    | \quad &\texttt{fun \textit{\rmfamily gname} [$x_1$, $x_2$] $x^{*}$ = $t_1$} \\
    & \texttt{\ \ | \textit{\rmfamily gname}\ \texttt{.}$xa$ \hspace{.7cm} $x^{*}$ = $t_2$} \\
    \mathcal{T} \ni t ::= & \texttt{\textit{\rmfamily fname} $x^{*}$} \\
    | \quad & \texttt{\textit{\rmfamily gname} $x^{*}$} \\
    | \quad & \texttt{if $ea_1 = ea_2$ then $t_1$ else $t_2$} \\
    \mathcal{E} \ni e ::= & \texttt{[$e_1$, $e_2$]} \\
    | \quad & x \\
    | \quad & ea \\
    \mathcal{A} \ni ea ::= & \texttt{'}s \\
    | \quad & \texttt{.}xa
  \end{align*}
  \caption{Syntax}
\label{fig:syntax}
\end{figure}


\section{Operational Semantics}
Judgement form: $\sigma |-_\Gamma t => v$


\begin{figure*}\centering
  
  \begin{equation}
    \nfrac{
      ea_1/ \sigma = ea_2/\sigma \quad
      \sigma |-_\Gamma t_1 => v_1
    }{
      \sigma |-_\Gamma \texttt{if $ea_1$ = $ea_2$ then $t_1$ else $t_2$} => v_1
    } \tagsc{If-True}
  \end{equation}

\begin{equation}
  \nfrac{
    ea_1/ \sigma \neq ea_2/\sigma \quad
    \sigma |-_\Gamma t_2 => v_2
  }{
    \sigma |-_\Gamma \texttt{if $ea_1$ = $ea_2$ then $t_1$ else $t_2$} => v_2
  } \tagsc{If-False}
\end{equation}

\begin{equation}
  \nfrac{
    \begin{array}{c}
      \Gamma(\textit{gname}) =
      \begin{array}{l}
        \texttt{fun \textit{\rmfamily gname} [$x_1$, $x_2$] $x_3$ $\ldots$ $x_n$ = $t_1$} \\
        \texttt{\ \ | \textit{\rmfamily gname}\ \texttt{.}$xa$ \hspace{.7cm} $x_3'$ $\ldots$ $x_n'$ = $t_2$} \\
      \end{array} \\
      e_1/\sigma = \texttt{[$d_1$, $d_2$]} \\
      \sigma\{x_1 |-> d_1, x_2 |-> d_2, x_3 |-> e_2/\sigma, \ldots, x_n |-> e_n/\sigma\} |-_\Gamma t_1 => v_1
    \end{array}
  }{
    \sigma |-_\Gamma \texttt{\textit{\rmfamily gname} $e_1$ $e_2$ $\ldots$ $e_n$} => v_1
  } \tagsc{Call-G-Cons}
\end{equation}

\begin{equation}
  \nfrac{
    \begin{array}{c}
      \Gamma(\textit{gname}) =
      \begin{array}{l}
        \texttt{fun \textit{\rmfamily gname} [$x_1$, $x_2$] $x_3$ $\ldots$ $x_n$ = $t_1$} \\
        \texttt{\ \ | \textit{\rmfamily gname}\ \texttt{.}$xa$ \hspace{.7cm} $x_3'$ $\ldots$ $x_n'$ = $t_2$} \\
      \end{array} \\
      e_1/\sigma = \texttt{'$s$} \\
      \sigma\{xa |-> \texttt{'}s, x_3' |-> e_2/\sigma, \ldots, x_n' |-> e_n/\sigma\} |-_\Gamma t_2 => v_2
    \end{array}
  }{
    \sigma |-_\Gamma \texttt{\textit{\rmfamily gname} $e_1$ $e_2$ $\ldots$ $e_n$} => v_2
  } \tagsc{Call-G-Atom}
\end{equation}

\begin{equation}
  \nfrac{
    \begin{array}{c}
      \Gamma(\textit{fname}) =
        \texttt{fun \textit{\rmfamily fname} [$x_1$, $x_2$] $x_3$ $\ldots$ $x_n$ = $t_1$}
        \\
      \sigma\{x_1 |-> e_1/\sigma, x_2 |-> e_2/\sigma, \ldots, x_n |-> e_n/\sigma\} |-_\Gamma t_2 => v_2
    \end{array}
  }{
    \sigma |-_\Gamma \texttt{\textit{\rmfamily fname} $e_1$ $e_2$ $\ldots$ $e_n$} => v_2
  } \tagsc{Call-G-Atom}
\end{equation}


\caption{Operational Semantics for }
\label{fig:semantics}
\end{figure*}

\bibliographystyle{abbrvnat}
\bibliography{../bibliography}

% \appendix
% \section{Appendix Title}




\end{document}
